\section{Related Work}

\subsection{Traditional Battery Thermal Modeling}
Traditional battery thermal models are broadly categorized into numerical models and lumped-parameter models. Numerical methods, such as Computational Fluid Dynamics (CFD) and Finite Element Methods (FEM), solve discretized three-dimensional heat conduction and convection equations \cite{zhang2017surrogate, ouyang2022surrogate}. Although these models provide spatial temperature profiles, their computational demand precludes real-time execution in vehicle microcontrollers. 

To enable real-time operation, lumped-parameter models are deployed. Single-node lumped models assume uniform cell temperature, modeling heat transfer via a single thermal capacity and convective resistance. For cylindrical or large-format pouch cells with significant internal temperature gradients, dual-node models (representing cell core and casing separately) are used \cite{saha2014lumped}. Equivalent circuit models (ECMs) map thermal dynamics to electrical parameters \cite{kaliaperumal2024thermal, xie2021enhanced, li2021novel}. While fast, lumped models require parameter calibration and struggle to resolve core-surface transient phase lags under dynamic drive cycles.

\subsection{Data-Driven and Digital Twin Models}
With increased sensor availability, machine learning models have been applied to battery state prediction. Multi-Layer Perceptrons (MLPs), Recurrent Neural Networks (RNNs), and Transformers have been explored for thermal estimation and degradation forecasting \cite{saravanan2022machine}. These data-driven models form core elements of battery digital twins for predictive monitoring \cite{lyu2024digitaltwin, patil2024review, semeraro2024battery, li2021review}. 

However, purely data-driven black-box models do not enforce physical conservation laws, making them susceptible to predictions inconsistent with thermodynamics under out-of-distribution operating conditions \cite{li2021review}. In this work, we focus on feedforward architectures as the primary baseline to isolate physics loss regularization without confounding state update mechanisms, acknowledging that recurrent and continuous-time neural ODE architectures represent valuable directions for future benchmarking.

\subsection{Physics-Informed Neural Networks (PINNs)}
Physics-Informed Neural Networks incorporate differential equations as soft regularization terms during training \cite{raissi2019physicsinformed}. By penalizing residuals of physical laws, PINNs constrain the parameter space and promote predictions consistent with the governing model.

In battery systems, PINNs have been investigated for state-of-charge (SOC) and state-of-health (SOH) estimation using electrochemical equations \cite{li2021physics, li2023hybrid, pollo2024coupling, ke2024physics}. For thermal estimation, PINNs have been applied to model spatial thermal distributions and anomaly detection \cite{chen2022physics, wang2024physicsinformed, ji2024physics, vairo2024physics, paoletti2024multip, ji2024realtime}. In this manuscript, PINN thermal ODE loss formulation is combined with trainable parameter estimation rather than treated as an isolated contribution.

\subsection{Exposure Bias and Dynamic Rollout}
In dynamic time-series forecasting, models are frequently trained using single-step teacher forcing \cite{ranzato2015sequence}. During inference, the model operates recursively, feeding prior predictions as inputs. This objective mismatch, termed exposure bias, leads to trajectory drift over extended rollouts \cite{ranzato2015sequence}.

While exposure bias has been studied in language and sequence modeling \cite{bengio2015scheduled}, its effect on multi-step physical system rollouts warrants further study. In high-frequency thermal digital twin rollouts ($dt = 0.01$~s corresponding to $150,000$ steps per cycle), sub-millikelvin single-step errors compound over time. We evaluate input noise injection during training as a practical mechanism to expose the network to out-of-trajectory states, improving multi-step rollout stability. The novelty of our approach lies in the \textbf{hybrid combination} of this exposure-bias mitigation with physics-informed regularization and parameter estimation within a battery digital twin context.
